% Иллюстрация к теме «внутренняя энергия идеального газа»:
% степени свободы двухатомной молекулы — поступательное движение центра масс
% и вращение вокруг двух взаимно перпендикулярных осей (y и z).
% Компиляция: pdflatex diatomic-molecule-figure.tex
\documentclass[tikz,border=8pt]{standalone}
\usepackage[T2A]{fontenc}
\usepackage[utf8]{inputenc}
\usepackage[russian]{babel}
\usetikzlibrary{arrows.meta,calc}

\begin{document}
\begin{tikzpicture}[
    >={Stealth[length=2.4mm]},
    % --- стили в духе первого рисунка ---
    axis/.style = {->, semithick},
    vel/.style  = {->, very thick, blue!55!black},
    rot/.style  = {->, line width=1pt, red!70!black},
    note/.style = {gray!30!black, font=\small},
  ]

  % ================= оси =================
  % молекула лежит вдоль оси x
  \draw[axis] (-2.5,0) -- (3.3,0) node[below right=-2pt] {$x$};
  \draw[axis] (0,-0.7) -- (0,2.95) node[above] {$y$};
  \draw[axis] (0,0) -- (-2.36,-1.50) node[below left=-2pt] {$z$};

  % ================= вращение вокруг оси y =================
  % эллипс с разрывом сверху (там ось проходит «перед» дугой)
  \draw[rot] ($(0,2.05)+(100:{0.8} and {0.27})$)
      arc[start angle=100, end angle=440, x radius=0.8, y radius=0.27];
  \node[red!70!black] at (1.05,2.05) [right] {$\omega_y$};

  % ================= вращение вокруг оси z =================
  \begin{scope}[rotate around={122.5:(-1.60,-1.02)}]
    \draw[rot] ($(-1.60,-1.02)+(110:{0.78} and {0.26})$)
        arc[start angle=110, end angle=440, x radius=0.78, y radius=0.26];
  \end{scope}
  \node[red!70!black] at (-1.02,-1.55) [right] {$\omega_z$};

  % ================= молекула =================
  % связь
  \draw[line width=4.5pt, gray!55, line cap=round] (-1.0,0) -- (1.0,0);
  % атомы
  \shade[ball color=gray!45] (-1.1,0) circle (0.45);
  \draw[gray!60, thin] (-1.1,0) circle (0.45);
  \shade[ball color=gray!45] (1.1,0) circle (0.45);
  \draw[gray!60, thin] (1.1,0) circle (0.45);

  % ================= центр масс и его скорость =================
  \fill (0,0) circle (0.06);
  \draw[gray!70, thin] (0.62,-0.78) -- (0.07,-0.09);
  \node[note, anchor=west] at (0.58,-0.85) {центр масс};

  \draw[vel] (0,0) -- (1.05,1.35) node[above right=1pt] {$\vec v$};

\end{tikzpicture}
\end{document}
